\chapter{Amenazas a la validez}
\label{cap:amenazas-validez}

Este cap\'itulo documenta, sin minimizarlas, las limitaciones metodol\'ogicas
que afectan la interpretaci\'on de los resultados presentados en los
cap\'itulos~\ref{cap:resultados-jaime}--\ref{cap:resultados-zaida}. Ninguna
medici\'on de este informe se presenta como libre de amenazas a su
validez.

\section{Validez interna}

\begin{itemize}[nosep]
  \item Todas las mediciones de seguridad y rendimiento se ejecutaron en
    \textbf{una sola m\'aquina de desarrollo} y contra \textbf{una \'unica
    instancia} del backend; no se control\'o la interferencia de otros
    procesos del sistema operativo durante las corridas.
  \item El certificado TLS es \textbf{autofirmado}, generado localmente
    para esta evaluaci\'on; el protocolo y el cifrado negociados
    (TLSv1.3, \texttt{TLS\_AES\_256\_GCM\_SHA384}) son reales, pero el
    canal de confianza del certificado en s\'i no lo es.
  \item El hit ratio de cach\'e se midi\'o con dos m\'etodos que producen
    cifras distintas ($\approx$49{,}98\,\% global frente a $\approx$100\,\%
    aislado); la elecci\'on del m\'etodo aislado como el m\'as
    representativo es una decisi\'on metodol\'ogica del equipo, documentada
    expl\'icitamente para que un lector externo pueda juzgarla.
  \item El rate limiting de login se evalu\'o \'unicamente en su
    implementaci\'on en memoria, por instancia \'unica; su comportamiento
    bajo m\'ultiples instancias del backend no fue evaluado porque no est\'a
    implementado (no es distribuido).
\end{itemize}

\section{Validez externa}

\begin{itemize}[nosep]
  \item Todas las pruebas se ejecutaron en un \textbf{entorno acad\'emico
    local}, no en un entorno de producci\'on ni de \emph{staging}
    equivalente a producci\'on; los resultados de rendimiento y seguridad
    no deben extrapolarse directamente a un despliegue en la nube con
    m\'ultiples usuarios concurrentes reales.
  \item Los usuarios de prueba (administrador sembrado, cuentas de
    registro) son \textbf{cuentas acad\'emicas de desarrollo}, no usuarios
    reales de una cl\'inica veterinaria; el comportamiento observado del
    sistema frente a datos de producci\'on reales (volumen, variedad,
    calidad de los datos) no fue evaluado.
  \item Lighthouse, cuando se ejecute, depender\'a del entorno de ejecuci\'on
    (hardware, red, versi\'on del navegador) en que se corra: sus
    puntuaciones no son un valor absoluto reproducible en cualquier
    m\'aquina, sino una medici\'on dependiente del entorno.
  \item Un posible \textbf{sesgo de participantes} en la prueba SUS
    (cuando se ejecute) es previsible si los diez participantes m\'inimos
    provienen mayoritariamente del entorno acad\'emico del equipo
    (compa\~neros, familiares con perfil t\'ecnico), en lugar de usuarios
    finales representativos de una cl\'inica veterinaria real.
\end{itemize}

\section{Validez de constructo}

\begin{itemize}[nosep]
  \item Las pruebas de integraci\'on que ejercitan la funci\'on nativa
    \texttt{fn_resumen_mascotas_por_especie}
    (\texttt{ResumenEspeciesIntegrationTest}) usan Testcontainers con
    PostgreSQL real; sin embargo, otras pruebas del backend usan H2 como
    base de datos en memoria bajo el perfil de pruebas, cuyo comportamiento
    \textbf{no es id\'entico} al de PostgreSQL real para funciones nativas
    PL/pgSQL --- por eso el resumen por especie se prob\'o expl\'icitamente
    con Testcontainers en vez de H2.
  \item \texttt{TokenBlacklistService} se prueba siempre con
    \texttt{@MockBean} en las pruebas unitarias/MockMvc; su comportamiento
    real contra una instancia de Redis en condiciones de fallo (ca\'ida de
    red, memoria llena) no se re-verific\'o en esta entrega.
  \item La ausencia de un SIEM centralizado significa que la
    ``auditor\'ia'' medida en este informe es la \textbf{emisi\'on
    correcta} de las l\'ineas de log \texttt{AUTH_AUDIT}, no su
    recolecci\'on, alertamiento o retenci\'on a largo plazo, que son
    aspectos distintos de un programa de auditor\'ia de seguridad.
\end{itemize}

\section{Validez de conclusi\'on}

\begin{itemize}[nosep]
  \item El tama\~no de muestra de cada corrida de k6 ($\sim$3\,200
    peticiones) es adecuado para el c\'alculo de percentiles e intervalos
    de confianza reportados, pero corresponde a una \'unica ventana de
    tiempo corta ($\sim$30--35~s); no se midieron tendencias a lo largo de
    per\'iodos prolongados (horas o d\'ias).
  \item Las 109 pruebas automatizadas del backend y la cobertura JaCoCo
    miden \textbf{cobertura de ejecuci\'on}, no ausencia de defectos: una
    l\'inea cubierta por una prueba no implica que est\'e libre de errores
    l\'ogicos no contemplados por esa prueba.
  \item Algunas evidencias (verificaci\'on manual de privilegios con
    \texttt{psql}, ejecuci\'on manual de \texttt{fn_resumen_mascotas_por_especie})
    se realizaron de forma manual y puntual, no como parte de un pipeline
    de integraci\'on continua que las reejecute autom\'aticamente en cada
    cambio.
\end{itemize}
